\section{The base rows and the theta square on a single linked bundle}

% BEGIN B74R1
\paragraph{B74R1. A single witness for the dependent data.}
Fix a parameter $p$, an original weighted point $q$ and a bundle $b$ at
that point, and suppose $w:\mathrm{BWit}(q,b)$. Put $u=w.x$,
$R=u.R$, $F=u.F$, $Y=u.Y$, $Z=u.Z$, $X=R.\mathrm{toCx3}(F)$.
The proof field of $w.y$ is a witness of $\mathrm{Holds}(R,F)$.
The equality in that same record identifies two elements of
$\mathrm{RPk}(b.R,b.S)$:
\[
 (b.P,b.C.c.\mathrm{face},b.C.\mathrm{Sp},\mathrm{zsc}_b,b.B)
 =
 (P_u,F.\mathrm{Dc},F.\mathrm{Sp},F.\mathrm{Zsc},
            \mathrm{realTBOf}(w.y.y)).
\]
Here $P_u$ is the row-side family built from $u$, and $\mathrm{zsc}_b$
uses the original pair of integers, scale and row-side family of $b$.
Any proposition depending on these five entries can be proved on the
right and transported to the left by substitution. The later entries
depend on the earlier ones, so projection to the first coordinate
alone would lose information. In the right-hand description the
realization is $R$, its face is $F$, and its chosen subspaces and
base rows are $X.\mathrm{zs}$ and $X.\mathrm{data}(F.\mathrm{cr},
\mathrm{false})$. All four identifications follow by expanding the
construction. The weak characterization comes from $u.\mathrm{pins}$,
whose phase field gives $R.\mathrm{phs}(j)=1$ for every lift $j$.
If one starts with the selected argument $w.a.x$, use $w.ax$ first.
This orders the identifications without choosing new column
spaces, bases or a second tuple reading.
% END B74R1

% BEGIN B74R2
\paragraph{B74R2. Base tests and smooth extensions for the stored map.}
Use B74R1 and write $K=D.\mathrm{what}_2.\mathrm{LL}$ and
$\sigma=A_3.A.v_0.\mathrm{embedding}$. For every lift $j$ there
are finitely many $\psi_a\in X.W(j)^\vee$ with
\[
 \begin{split}
 &\bigcap_a \ker\psi_a=0,\\
 &\psi_a\bigl(X.\theta'_{jt}(X.jL_{jt}v)\bigr)\in\sigma(K)
               \quad(t\text{ arbitrary},\ v\in A.FL(j,t)),\\
 &\psi_a\bigl(X.w0C_{0j}(e_k)\bigr)\in\sigma(K)
               \quad(1\le k\le X.dW(0,j)).
 \end{split}
\]
In the last line $0$ denotes the base level and $e_k$ runs through
its stored coordinate vectors. Indeed, take
$w.y.\mathrm{intersectionNumber\_baseChange}.\mathrm{prim}$ and
expand its $\mathrm{liftTestB}$ field. The quantifier for $\psi$
precedes both rationality clauses. This proves separation on the
entire target, with the tests in the twisted model used in BTH8.
It retains the distinction between $\sigma$ and the distinguished
embedding used by the semilinear geometric reading.

The same $\mathrm{Holds}$ has an auxiliary field
$R.\mathrm{toRealDat}.\mathrm{ExAt}$. Its $\mathrm{thF}$ entry gives,
for every $j,t$, a map satisfying $\mathrm{PThF}$. Apply the first
implication of the weak choice rule in $u.\mathrm{pins}.\mathrm{char}$
to that witness. The \emph{stored} $U=R.\mathrm{thF}(j,t)$ then
takes every vector into the smooth cuspidal multiplicity space
$\mathrm{Wall}_j$ and restricts on pure tensors of lowest type to
the prescribed integral on $\operatorname{im}(R.\mathrm{extV}(j,t))$.
In particular evaluating $U$ on one basis at the base tuple gives the
finite extensions of BTH4. The analogous $\mathrm{thG}$ implication
identifies the stored projected map via $e_W\,\mathrm{thG}
=\mathrm{postOn}(eWf)\,U$. These deductions use the given
$\mathrm{Holds}$ only downstream of its construction.
% END B74R2

% BEGIN B74R3
\paragraph{B74R3. Constructing the eight base row fields.}
On the realization of B74R1 put $H=w.y.\mathrm{intersectionNumber\_baseChange}$.
Use $s=1$, $H.\mathrm{ok}$, $H.\mathrm{sep}$ and
$h=H.\mathrm{prim}$. We describe the base-level construction
$\mathrm{primRows3}$ with underlying data $X.\mathrm{data}(F.\mathrm{cr},0)$.
The weak model gives the rational forms and the row and lift isometries.
From $h.\mathrm{hecke}$ use the chosen row and lift operators together
with their intertwining equations. Use these very lift operators in
$h.\mathrm{liftOnto}$. The row and lift tests give rationality of the
good subspaces via rational orthogonal complements. Apply the native
unitary eigensystem decomposition with these subspaces, the rational
adjoints, lower-level type, unitary equivariance, the decomposition
into lift parts and the preceding surjectivity. Both outputs use one
eigensystem choice. They give the base types (using separation) and
the base proviso (using the given Hecke conditions).
The primitive positivity supplies the positive level, and the tuple
tie supplies the height identity $H0$.

For the frames, the row model gives the stated row frame, the lift
frame has invertible Gram matrix by the rational form, and the
integral lattice assertion is $h.\mathrm{mLat}$. The source's
base lift identity identifies the complexification of the rational
lift matrix with the theta map on the cut coordinates. Its rational
entries follow from the base lift isomorphism onto its image, the
principal frame and the tests of B74R2. The tests embed the target
in a coordinate space. The test matrix on the independent frame
columns has an invertible maximal minor over the embedded field.
Solving in those rows expresses each tested lift column over that
field. This is the native base rationality argument before forming
the row record. On the row side
the base row matrix is the identity, so its image of each type
subspace is precisely its prescribed rational part. On the lift side
the type membership test and the lift identity give the forward
inclusion. Lift-matrix surjectivity gives a preimage of each vector of
the target part, and the reverse membership test puts that preimage
in the specified source part. This gives both images in the field
called $\mathrm{l1a}$. The row and lift matrix bijections, scalar
positivities and the row and lift scalar transformation laws are the
entries of $\mathrm{l1c}$.

Here is the dimension argument for the last field. For each type $p$
let $a_p$ be the dimension over $K$ of its base lift coordinate
subspace, $b_p$ that of the corresponding rational target part, and
$c_p$ the complex dimension of the analytic part intersected with
the image of the base frame. The rational lift matrix embeds the
first subspace in the second. An independent family in rational
coordinates remains independent after scalar extension, and the
complex frame is injective. Hence $a_p\le b_p\le c_p$.
The complex intersections are independent by orthogonality, so
$\sum c_p$ is bounded by the dimension of the image of the frame.
The base lift is injective with precisely that image after extension.
Its source dimension is at most $\sum a_p$ since the type spaces
span the base coordinates. Thus equality holds separately for every
$p$, including a zero-dimensional or missing part. The row version
uses the base row image and its orthogonal parts by the same dimension
bounds. The primitive complex equality $h.\mathrm{dimParts}$ now
identifies the two rational part dimensions. This is $\mathrm{l1d}$.

Together positivity, proviso, types, $H0$, frames, the two image
equalities, the basis and scalar identities, and the dimension
equalities are precisely the eight additional fields of
$\mathrm{IndefRows}$. Its underlying data is the specified $X.\mathrm{data}$
by construction. Weak and primitive data also give stability and
the in-block assertion. Transport stability, block membership and
the row witness along the realization's equality for its subspaces,
then use its dependent equality for $S_0$. With $\langle 1,H.\mathrm{ok}\rangle$
and $H.\mathrm{sep}$ these are the five entries of
$\mathrm{RowsHold0}_3(b.P)$.
% END B74R3

% BEGIN B74R4
\paragraph{B74R4. The square and rank on every original linked selection.}
The same $Y$ contains $\mathrm{RecTied}$. Its lift field evaluates
the geometric theta map on a used tuple whenever a multiplicity
vector realizes the prescribed lowest-type integral and its good
projection belongs to the target. Here $Y.\mathrm{tup}(t_0)=\mu_0$
is used. The source and target Hilbert readings in $Y$, the coordinate
law $Z.\mathrm{jW\_read}$, and the fields
$w.y.r.\mathrm{mtRead},w.y.r.\mathrm{mjRead}$ all refer to these maps.
Furthermore $u.\mathrm{div}$ supplies the independence, spanning
and lattice identification of the target basis. These data belong to
the row-side family and realization identified in B74R1.
B74R2 supplies the smooth extensions and the simultaneous base
tests. BTH7 therefore evaluates the stored theta map on the full
level, BTH10 kills the orthogonal complement of the good source,
and BTH11 reads its restriction in the stored target coordinates.
Writing a coordinate vector as $E_j\Pi_jv+(v-E_j\Pi_jv)$ gives
BTH12's square
\[
 D.\theta P\bigl(\mu_0,j+1,D.jMt(\mu_0,j+1)(Y.eMt(j,t_0)v)\bigr)
       =Z.jW(j)(M_j\Pi_jv).
\]
Here $M_j$ is the base theta matrix, and $E_j,\Pi_j$ are the
inclusion and normalized projection for $b.P$. B74R3 supplies the
row witness to the base-rank argument SLR4--SLR6 on this very bundle.

The quantifiers give a stronger simultaneous conclusion on linked
selections. Fix any original class $\mathcal A$, raw selector $\sigma_b$
and row condition $E$ for which the source's
$\mathrm{IsSel}(Q,\mathrm{IsLinked}(p),E,\mathcal A,\sigma_b)$
holds. For each member $\mathcal F$, each $x>\mathcal F.x_1$,
and each index at $x$, its defining clause gives
$\mathrm{Built}$ for $b=\sigma_b(\mathcal F,x,i)$.
Eliminate the dependent equality in $\mathrm{Built}$, then the one
in $\mathrm{linkOf}$ defining $\mathrm{IsLinked}$. What remains
is precisely a weighted point $q$ and a witness $w$ for that
same $b$. Apply B74R1--B74R3 and the square above before transporting
the conjunction back. Thus all four BTH12 duties hold on this
selection at \emph{every} admissible index. The original eventual
rank statement follows with $x_0=\mathcal F.x_1+1$, independent of
the point and the lift. Taking $E$ identically true also proves the
source's universal selected-row implication ($\mathrm{Row.of\_sel}$).
It can be conjoined with the other row duties of DEX8.

The existence step in DEX8 is separate: it requires a linked bundle
at every required scale of every prescribed good point to obtain
the design class with its existence properties. Unfolding a given
$\mathrm{BWit}$ here uses the primitive rational tests and rows
already packed into its $\mathrm{Holds}$. To produce the original
selection one must construct those data upstream. This identifies
precisely where a new arithmetic or geometric existence proof enters.
% END B74R4
